\documentclass[11pt,a4paper]{article}

\usepackage[margin=2.6cm]{geometry}
\usepackage{amsmath,amssymb}
\usepackage{booktabs}
\usepackage{graphicx}
\usepackage{xcolor}
\usepackage{enumitem}
\usepackage[colorlinks=true,linkcolor=blue!60!black,citecolor=blue!60!black,urlcolor=blue!60!black]{hyperref}
\graphicspath{{figures/}}

% Shared notation with the local/global addendum.
\newcommand{\zlat}{\mathbf{z}}
\newcommand{\act}{\mathbf{a}}
\newcommand{\abar}{\bar{\mathbf{a}}}

% Figures produced on Drive by the results notebook; degrade gracefully when absent.
\newcommand{\drivefigure}[3]{%
  \IfFileExists{figures/#1}{\includegraphics[width=#2]{#1}}{%
    \fbox{\parbox[c][4.0cm][c]{0.9\linewidth}{\centering\itshape
      Figure pending: copy \texttt{#3} from the Drive \texttt{\detokenize{_summary}}
      directory into \texttt{\detokenize{reports/figures/#1}} and rebuild.}}}}

\title{DINO-WM on PointMaze under a Commodity-GPU Budget:\\
Latent-Cached Reproduction, Low-Data Transfer, and Planning Results\\
{\large Technical report --- method, implementation, and results}}
\author{Thomas Georges\thanks{Implementation assistance: Claude (Anthropic). Code:
\texttt{wm-prediction} repository, \texttt{src/wm\_poc/dino\_wm}. This report has
an addendum, \emph{Local/Global World-Model Planning with DINO-WM}
(\texttt{reports/local\_global\_report.tex}), which builds a local/global
planning track on the checkpoint and cache documented here.}}
\date{June 2026 --- v1.0 (50-evaluation results; 200-evaluation pass pending)}

\begin{document}
\maketitle

\begin{abstract}
We reproduce DINO-WM~\cite{dinowm2025} --- an action-conditioned latent
transition model trained over frozen DINOv2 patch features for zero-shot
visual planning --- on the PointMaze task, under the constraint that the full
experiment must run on freely available Colab hardware. The stock training
loop required more than 30 hours per epoch on this hardware, independent of
batch size; we identify the cause as per-sample episode reload compounded by
online re-encoding of the \emph{frozen} encoder, and replace it with a
precomputed-latent training path that reproduces the upstream data
distribution bit-for-bit while cutting per-step cost by ${\approx}21\times$
(58~minutes per epoch on a T4, an estimated 3--6 on an A100). On this
foundation we run three experiments: a full-data predictor (2000 rollouts), a
low-data scratch baseline (300 rollouts), and a low-data fine-tune initialized
from the full-data checkpoint. The fine-tuned predictor reaches the scratch
run's converged latent-prediction loss at its \emph{first} epoch
($\sim$30$\times$ fewer updates), and in closed-loop CEM planning in the real
MuJoCo environment it delivers $2.75\times$ the scratch baseline's success rate
($0.275$ vs $0.10$ at 200 evaluation episodes, $p<10^{-5}$) while
\emph{matching} the full-data model trained on $6.7\times$ more data ($0.275$
vs $0.23$, $p{=}0.30$). A striking dissociation emerges: the scratch and
fine-tune predictors land close in validation loss yet differ by $2.75\times$
in planning success, evidence that latent-prediction MSE is a weak proxy for
downstream control. All numbers are read from recorded planner telemetry (the
latest \texttt{final\_eval} per run, at 200 evaluation episodes). An addendum
(\S\ref{sec:localglobal}) reuses this checkpoint as the trusted global model for
a coupled local/global planning study over six planners: a global re-score gate
makes cheap first-order refinement \emph{safe}, recovering global-CEM-quality
plans ($0.97$ success) where an unguarded hybrid collapses to $0.54$, while the
purely local planners run $100$--$1000\times$ cheaper per episode.
\end{abstract}

\section{Introduction}
\label{sec:intro}

A world model is useful for control to the extent that planning against it
produces actions that succeed in the real environment. DINO-WM~\cite{dinowm2025}
makes a specific bet about how to obtain such a model cheaply: rather than
learning perception and dynamics jointly, it \emph{freezes} a powerful
self-supervised visual encoder (DINOv2~\cite{dinov2}) and learns only an
action-conditioned transition model over its patch features, then plans by
sampling-based optimization of latent goal distance. The appeal is data- and
compute-efficiency --- no pixel reconstruction is needed in the planning loop,
and the encoder is reused unchanged across tasks.

This report documents a faithful reproduction of that pipeline on PointMaze,
developed under a hard practical constraint: it must run end to end on the
hardware actually available to us (Colab, predominantly T4; A100 rarely).
That constraint turned out to be the technically interesting part. The stock
pipeline does not merely run slowly on a T4 --- it is structurally infeasible,
projecting more than 30 hours per epoch regardless of batch size. Diagnosing
and removing that cost \emph{without} altering the experiment is the
methodological contribution of \S\ref{sec:training}; it is what makes the
results of \S\ref{sec:results} obtainable at all.

\paragraph{Contributions.} (i) A latent-cache training path for DINO-WM that
removes the dominant data-path cost while provably preserving the upstream
training distribution (\S\ref{sec:training}); (ii) a faithful, fully resumable
reproduction of DINO-WM predictor training and CEM planning on PointMaze,
runnable on commodity GPUs (\S\ref{sec:implementation}); (iii) a low-data
transfer study --- scratch vs.\ fine-tune from a full-data checkpoint --- with
both latent-prediction and real-environment planning metrics
(\S\ref{sec:results}); and (iv) the observation that latent-prediction loss
and planning success dissociate sharply, which we argue is a methodological
warning for world-model evaluation (\S\ref{sec:discussion}). This reproduction
is also the foundation for the coupled local/global planning study of
\S\ref{sec:localglobal}, which reuses this checkpoint and cache unchanged as a
trusted, non-differentiated global model.

\section{Background}
\label{sec:background}

\paragraph{DINO-WM.} DINO-WM~\cite{dinowm2025} encodes each frame to frozen
DINOv2 patch tokens $\zlat_t \in \mathbb{R}^{P\times D}$ and trains a
transformer predictor to advance the latent state conditioned on actions,
with a latent-consistency objective and no decoder required for planning.
Planning is visual goal reaching: given current and goal observations,
sampling-based MPC minimizes the predicted latent distance to the goal
encoding. The model's defining property for our purposes is that the encoder
is frozen --- its outputs are constants of the dataset, a fact the training
loop should but does not exploit.

\paragraph{Self-supervised features for control.} The broader pattern of
planning on top of frozen self-supervised representations
(DINOv2~\cite{dinov2}; the JEPA world-model family~\cite{vjepa2}) trades
end-to-end representation learning for modularity and sample efficiency.
DINO-WM is a concrete, planning-focused instance.

\paragraph{Sampling-based latent planning.} PlaNet~\cite{planet} and
Dreamer~\cite{dreamer} plan or learn policies inside a learned recurrent
latent model; TD-MPC2~\cite{tdmpc2} pairs a learned latent model with
sampling-based MPC. DINO-WM keeps the model fixed (frozen encoder, separately
trained predictor) and plans with the cross-entropy method (CEM), which
requires only forward rollouts and is the planner reproduced here.

\section{Problem setting and notation}
\label{sec:setting}

PointMaze episodes are sequences of frames with raw 2-D actions
$\act_t \in \mathbb{R}^{2}$, bounds $[-1,1]^2$, and proprioception (agent
position/velocity). Following DINO-WM, one \emph{model step} advances
$f$ raw steps (frameskip $f{=}5$); the model-step action is the folded block
$\abar_k = [\act_{kf};\dots;\act_{(k+1)f-1}] \in \mathbb{R}^{fA}$. Each frame
is encoded to $\zlat \in \mathbb{R}^{P\times D}$ with $P{=}196$ patches and
$D{=}384$ (DINOv2 ViT-S/14 at $224$\,px, whose internal resize yields a
$14\times14$ patch grid). The predictor consumes $H{=}3$ context frames
(\texttt{num\_hist}) and predicts the next latent ($\texttt{num\_pred}{=}1$),
trained autoregressively over windows. Episode-level train/validation splits
are drawn from a seeded permutation of the uncapped episode count, so split
membership is invariant to how many rollouts a given run uses.

\section{The DINO-WM model as reproduced}
\label{sec:model}

\paragraph{Encoder (frozen).} \texttt{dinov2\_vits14}, feature key
\texttt{x\_norm\_patchtokens}, weights frozen throughout
(\texttt{train\_encoder=false}). Input frames are preprocessed with the
upstream transform (resize/centre-crop to $224$\,px, normalize), then resized
to the $196$-px encoder input that produces the $14{\times}14$ patch grid.

\paragraph{Predictor.} A ViT predictor (\texttt{models.vit.ViTPredictor};
depth 6, 16 heads, MLP width 2048, dropout 0.1) over the flattened
$H \times P = 3 \times 196 = 588$ token sequence. Actions and proprioception
are embedded ($\text{action\_emb}{=}\text{proprio\_emb}{=}10$) and, with
\texttt{concat\_dim}=1, concatenated along the feature axis (tiled across
patches), giving the predictor working dimension
$D + 10 + 10 = 404$ (observed as ``\texttt{Model emb\_dim: 404}'' in the run
logs). A causal mask is applied across the temporal axis.

\paragraph{Training objective.} With $\hat{\zlat}$ the predicted next latent
and $\zlat$ the (detached) target encoding,
\begin{equation}
\mathcal{L} \;=\; \big\lVert \hat{\zlat} - \mathrm{sg}[\zlat] \big\rVert_2^2 ,
\label{eq:loss}
\end{equation}
the latent-consistency MSE, reported per epoch as a training and a validation
component (the ``h-step'' validation loss). The runs in this report are
\emph{decoder-free} (\texttt{has\_decoder=false}): no pixel reconstruction is
trained, which both halves memory and matches the planning objective, since
planning never decodes. The action and proprioception encoders are trained
jointly with the predictor; the DINOv2 encoder is shared and frozen by
construction.

\section{Feasible training: the latent-cache path}
\label{sec:training}

\subsection{Why the stock loop costs 30 hours per epoch}

Profiling the upstream data path identified three compounding causes, none of
which batch size can address (confirmed empirically: batch 4 and batch 32 both
project ${\sim}19$--$21$\,h/epoch on a T4).

\begin{enumerate}[leftmargin=1.4em]
  \item \textbf{Per-sample episode reload.} Upstream
        \texttt{TrajSlicerDataset} fetches a 4-frame window by
        \texttt{torch.load}-ing the \emph{entire} episode image tensor
        ($\sim$93 frames) and transforming \emph{every} frame, then discarding
        all but four. At ${\sim}147$k windows/epoch this re-reads and
        re-transforms ${\sim}1600\times$ the dataset per epoch, through the
        Google Drive FUSE mount.
  \item \textbf{Online re-encoding of frozen features.} With
        \texttt{train\_encoder=false} the DINOv2 outputs are dataset
        constants, yet every frame of every window is re-encoded every epoch
        --- pure recomputation.
  \item \textbf{Operational confounders.} Single-threaded loading
        (\texttt{num\_workers=0}, an OOM-era setting) and reads from Drive
        rather than local disk.
\end{enumerate}

\subsection{Encode once, reuse forever}

Because the encoder is frozen, its outputs can be precomputed. Three
components, installed into the upstream checkout at launch time as
marker-guarded patches:

\begin{enumerate}[leftmargin=1.4em]
  \item \textbf{One-time precompute.} Each episode is read once, preprocessed
        \emph{exactly} as upstream does, encoded under autocast through the
        frozen encoder, and written as a per-episode \texttt{fp16} \texttt{.npy}
        plus a manifest (encoder identity, patch geometry, episode lengths,
        dataset coverage). The cache is ${\sim}30$\,GB for all 2200 episodes,
        held on session-local disk (a Drive-hosted cache would reinstate the
        FUSE bottleneck; the scripts warn if so configured), and the process
        is idempotent and resumable.
  \item \textbf{A drop-in latent dataset}, selected purely by Hydra override.
        It subclasses the upstream dataset, so actions, states,
        proprioception, normalization statistics, and the seeded train/val
        split are \emph{bit-identical} to the online path (verified directly
        against the upstream loader); its slicer memory-maps only the
        ${\sim}0.6$\,MB of latents each window needs instead of materializing a
        ${\sim}13$\,MB episode.
  \item \textbf{An input-dispatch bypass} in the model's \texttt{encode\_obs}:
        4-D latent inputs skip the encoder; 5-D image inputs take the original
        path. The real DINOv2 weights remain in the model and its checkpoints,
        so planning (which encodes real observations) is untouched.
\end{enumerate}

This is an \emph{exact} reproduction, not an approximation: the only quantities
that change are when and how often the frozen encoder runs. A guard rejects
combining the latent path with the (pixel-target) decoder at command-build
time.

\subsection{Throughput}

\begin{table}[h]
\centering
\small
\begin{tabular}{lllll}
\toprule
Configuration & Steps/epoch & Step time (T4) & Epoch (T4) & Epoch (A100, est.) \\
\midrule
Stock, batch 4              & 36{,}815 & $\sim$3\,s    & $>$30\,h          & --- \\
Stock, batch 32             & 4{,}602  & 15.9--17\,s   & 19--21\,h (proj.) & --- \\
Latent cache, batch 32, stride 1 & 4{,}602 & $\sim$0.76\,s & $\sim$58\,min & 6--12\,min \\
Latent cache, batch 32, stride 2 & 2{,}330 & 0.72--0.76\,s & 31\,min (meas.) & 3--6\,min \\
\bottomrule
\end{tabular}
\caption{Training throughput before and after the latent-cache path. The
residual step cost is the depth-6 ViT predictor itself ($\sim$2\,TFLOPs
fwd+bwd at batch 32 over 588 tokens) --- the quantity the experiment is meant
to measure. ``Stride 2'' subsamples window start positions (consecutive
stride-1 windows overlap by 19/20 frames) and was used only for the full run
on T4; the low-data pair uses the full stride-1 window set.}
\label{tab:throughput}
\end{table}

\section{Planning and evaluation}
\label{sec:planning}

Planning uses the upstream CEM~\cite{cem}/MPC path unchanged. Given a start and
a goal observation, CEM optimizes an open-loop action sequence to minimize the
predicted latent distance to the goal encoding:
\begin{equation}
J(\abar_{1:H}) \;=\; \tfrac{1}{PD}\big\lVert \hat{\zlat}_H(\abar_{1:H}) - \zlat_{\mathrm{goal}} \big\rVert_2^2 ,
\label{eq:objective}
\end{equation}
iterating a sample/score/refit loop (population 100, 10 iterations, goal
horizon $G{=}5$, \texttt{goal\_source=random\_state}). Crucially, while CEM
\emph{optimizes} an imagined latent quantity, \textbf{success is measured in
the real MuJoCo environment}: the planned actions are executed in the
simulator and success is the environment's own goal-reaching test. Reported
success rates are therefore real closed-loop control outcomes, not the model's
self-assessment. Each evaluation also records, per CEM iteration, the current
best plan's environment success, mean state distance to goal, and the
divergence between the imagined and executed latent trajectories (the model's
own rollout error).

\section{Implementation}
\label{sec:implementation}

The pipeline follows a two-notebook convention: an operational notebook (02)
for setup, precompute, training, fine-tuning, and planning --- every heavy step
delegated to a gated CLI script against a \emph{pristine} upstream checkout ---
and a read-only results notebook (07) that aggregates artifacts into tables,
curves, and embedded planning videos. Reproducibility and reliability features:

\begin{itemize}[leftmargin=1.4em]
  \item \textbf{Pristine upstream, runtime patches.} All behavioral changes
        (mixed precision, latent cache, validation \texttt{no\_grad},
        fine-tune initialization, decoder-free video recording) are
        marker-guarded, idempotent, reversible textual patches applied at
        launch, with backups; the upstream repository is never vendored or
        edited in place. Patch application and composition are unit-tested
        against the real upstream source.
  \item \textbf{Single experiment, hardware-only profiles.} GPU-tier configs
        differ only in throughput knobs (precision, dataloader workers,
        wall-clock caps, and --- on T4 only --- window stride); model sizes,
        optimizer budgets, episode counts, and run names are shared so a run
        continues across tiers.
  \item \textbf{Resumability.} Full epoch checkpoints plus a rolling
        intra-epoch state-dict checkpoint (every 2000 steps on A100, 500 on
        T4) with deterministic mid-epoch resume bound a disconnection to
        minutes of lost work.
  \item \textbf{Self-gating cells.} Each launch cell checks durable state
        before spending GPU time: completed training is skipped, partial
        training resumes, and planning re-triggers only when the recorded
        evaluation does not match the configured episode count --- making a
        top-to-bottom run idempotent.
  \item \textbf{Smoke crash test.} A 20-rollout configuration exercises the
        exact production pipeline (patches, precompute, latent training,
        checkpointing, a 2-episode planning pass) in $\sim$10--15 minutes; it
        caught a planner/latent-dataset incompatibility, a checkpoint
        resolution bug, and a blank-checkpoint crash before any cost full-run
        time.
  \item \textbf{Testing.} ${\sim}95$ unit tests (runnable without torch) cover
        patch application/composition, command rendering, run-state gates, and
        metrics extraction; latent-dataset parity against the upstream loader
        is validated separately with torch.
\end{itemize}

\section{Experimental setup}
\label{sec:setup}

\begin{table}[h]
\centering
\small
\resizebox{\textwidth}{!}{%
\begin{tabular}{@{}llll@{}}
\toprule
\textbf{Component} & \textbf{Setting} & \textbf{Component} & \textbf{Setting} \\
\midrule
Encoder & DINOv2 ViT-S/14, frozen, $224$\,px & Predictor & ViT, depth 6, 16 heads, MLP 2048 \\
Latents & $P{=}196$, $D{=}384$, fp16 cache & Working dim & $404$ ($384{+}10{+}10$, concat) \\
Context / pred & $H{=}3$, \texttt{num\_pred}{=}1 & Frameskip & $f{=}5$ \\
Batch size & 32 & Predictor lr & $5{\times}10^{-5}$ (scratch) \\
Action-enc.\ lr & $5{\times}10^{-4}$ (scratch) & Fine-tune lr & $10^{-5}$ / $10^{-4}$ \\
CEM pop.\ / iters & $100$ / $10$ & Goal horizon $G$ & $5$ \\
Goal source & \texttt{random\_state} & Planning evals & $200$ \\
\bottomrule
\end{tabular}}
\caption{Configuration. Precision is bf16 on A100/L4 and fp16 on T4 (Turing
emulates bf16 slowly); all runs are seed 0.}
\label{tab:config}
\end{table}

Three runs enter the results (Table~\ref{tab:runs}). All train the predictor
only, on cached latents, with identical seeds, normalization, and splits.

\begin{table}[h]
\centering
\small
\begin{tabular}{@{}lllll@{}}
\toprule
Run & Data (train+val) & Epochs & Initialization & Planning evals \\
\midrule
Full no-decoder    & 2000 + 200 & 10 & random          & 200 \\
Low-data scratch   & 300 + 60   & 30 & random          & 200 \\
Low-data fine-tune & 300 + 60   & 20 & full checkpoint & 200 \\
\bottomrule
\end{tabular}
\caption{Runs. Fine-tuning loads the predictor and action/proprio encoders
from the full run's checkpoint and continues at reduced learning rates; the
frozen DINOv2 encoder is shared. The full run was trained on a T4 with
stride-2 window subsampling (Table~\ref{tab:throughput}).}
\label{tab:runs}
\end{table}

\section{Results}
\label{sec:results}

All numbers below are read directly from recorded artifacts --- the
\emph{latest} \texttt{final\_eval} row of each run's planner
\texttt{logs.json}, at 200 evaluation episodes per checkpoint (Wilson
half-width ${\approx}0.06$).

\subsection{Training loss: fine-tuning starts where scratch ends}
\label{sec:results-loss}

\begin{figure}[h]
\centering
\drivefigure{pointmaze_training_curves.png}{0.78\linewidth}{pointmaze\_training\_curves.png}
\caption{Per-epoch training (dashed) and validation (solid) latent-prediction
loss. The low-data scratch run descends from $\approx 2.3$ to $\approx 0.02$
over 30 epochs; the fine-tune begins at $\approx 0.015$ at epoch 1 --- already
the scratch run's converged level --- and stays there for all 20 epochs.}
\label{fig:curves}
\end{figure}

The fine-tuned predictor attains at \emph{epoch 1} the validation loss the
scratch baseline reaches only after ${\sim}30$ epochs: pre-training delivers
the converged predictor for ${\sim}30\times$ fewer updates on the low-data
task. The honest caveat is that scratch \emph{does} converge to the same
validation loss by epoch 30 (Figure~\ref{fig:scratchft}). On the
latent-prediction objective alone, the two models therefore converge to the
same order of magnitude --- which makes the planning result that
follows all the more pointed.

\begin{figure}[h]
\centering
\drivefigure{pointmaze_scratch_vs_finetune.png}{0.66\linewidth}{pointmaze\_scratch\_vs\_finetune.png}
\caption{Final validation h-step loss, scratch vs.\ fine-tune: $\approx 0.021$
vs $\approx 0.012$ --- the same order of magnitude, in sharp contrast to their
$2.75\times$ gap in planning success (Figure~\ref{fig:success}).}
\label{fig:scratchft}
\end{figure}

\subsection{Planning success in the real environment}
\label{sec:results-success}

\begin{figure}[h]
\centering
\includegraphics[width=0.72\linewidth]{success_rates.pdf}
\caption{Closed-loop CEM planning success in MuJoCo PointMaze, 200 evaluation
episodes per checkpoint, with 95\% Wilson intervals. Exact values from the
planner logs: fine-tune $0.275$, full-data $0.23$, scratch $0.10$.}
\label{fig:success}
\end{figure}

At $n{=}200$ the Wilson intervals tighten to $\pm 0.06$, and two clean results
stand out. \textbf{Fine-tune $\boldsymbol{>}$ scratch} is decisive ($0.275$ vs
$0.10$; two-proportion $z{=}4.5$, $p<10^{-5}$): pre-training on the full corpus
and adapting on 300 rollouts yields $2.75\times$ the planning success of
training on those same 300 rollouts from scratch. \textbf{Fine-tune vs.\
full-data} ($0.275$ vs $0.23$) is statistically indistinguishable
($p{=}0.30$): fine-tuning \emph{matches} a model trained on $6.7\times$ more
data. The full-data model's $0.23$ is moreover a lower bound: it was trained
on a T4 with stride-2 window subsampling (half the gradient updates per epoch
of the paper-exact stride-1 set), so an A100 stride-1 training of the identical
configuration should only improve it --- which would make the fine-tune's
parity with it the more striking. (The larger sample also moved all three
estimates down from the noisier 50-eval interim, $0.34/0.26/0.12 \to
0.275/0.23/0.10$, without disturbing the ranking or either conclusion.)

\subsection{The loss/planning dissociation}

The central empirical finding is the contrast between
Figures~\ref{fig:scratchft} and~\ref{fig:success}: scratch and fine-tune
converge to the same order of magnitude in validation latent-prediction loss
($\approx 0.021$ vs $\approx 0.012$) yet differ $2.75\times$ in
real-environment planning success ($0.10$ vs $0.275$, $p<10^{-5}$). The
${\sim}2\times$ residual loss gap is far too small to explain a near-tripling
of control success: latent-MSE convergence does not imply equal
control quality --- the fine-tuned model's latent geometry is evidently more
\emph{plannable} (smoother or better-calibrated for CEM's optimization)
despite matching the scratch model's average prediction error. The practical
consequence for world-model evaluation is direct: a planning metric must be
carried alongside the loss curve; the loss alone would have reported these two
models as equivalent.

\subsection{Planner diagnostics: CEM is genuinely optimizing}
\label{sec:results-cem}

\begin{figure}[h]
\centering
\begin{minipage}{0.49\linewidth}\includegraphics[width=\linewidth]{cem_success.pdf}\end{minipage}\hfill
\begin{minipage}{0.49\linewidth}\includegraphics[width=\linewidth]{cem_state_dist.pdf}\end{minipage}
\caption{Planner-internal telemetry per CEM iteration (exact values from
\texttt{logs.json}; for relaunched evaluations the latest attempt per
iteration is shown). Left: environment success of the current best plan rises
across the 10 iterations (full $0.135\to0.23$; fine-tune $0.135\to0.275$;
scratch $0.03\to0.10$). Right: mean state distance to goal falls
correspondingly (e.g.\ fine-tune $1.45\to0.93$).}
\label{fig:cem}
\end{figure}

\begin{figure}[h]
\centering
\includegraphics[width=0.58\linewidth]{cem_embed_div.pdf}
\caption{Divergence between imagined and executed visual embeddings per CEM
iteration --- the world model's own rollout error along the optimized plan.
The fine-tune and full models keep imagination closest to reality at the end
of optimization.}
\label{fig:embdiv}
\end{figure}

The iteration curves confirm CEM is exploiting the learned dynamics rather than
stumbling onto goals: success rises and goal distance falls monotonically over
the 10 CEM iterations for every trained model, with the scratch model
improving least --- consistent with its weaker, less plannable dynamics. The
embedding-divergence curve (Figure~\ref{fig:embdiv}) shows the better models
also keep imagination closer to executed reality, which is exactly the
property a planner needs.

\subsection{Planning videos}

Each evaluation records per-episode MP4s
(\texttt{output\_final\_<idx>\_<success$|$failure>.mp4}; 26 frames at 12\,fps,
the $5{\times}5{+}1$-step horizon) of the executed plan in the environment
against the goal image. (A patch was required for this: upstream records
videos only for decoder-bearing models, so the decoder-free checkpoints
produced none until the recorder was extended with a no-decoder branch ---
executed rollout vs.\ goal, imagination panel blacked out. Videos predating the
patch come only from a 1-epoch decoder-bearing smoke model and are not
representative.)

\section{Discussion and limitations}
\label{sec:discussion}

\paragraph{What these results establish.} A faithful DINO-WM PointMaze
reproduction is feasible on commodity hardware once the frozen-encoder cost is
removed, and on it, transfer from a full-data checkpoint yields a real,
measurable planning advantage on a low-data task --- with the efficiency
bonus that the fine-tuned predictor is converged after a single epoch. The
loss/planning dissociation is, to us, the most scientifically interesting
output and the firmest methodological lesson.

\paragraph{Threats to validity.} (i)~\emph{Sample size}: 200 planning episodes
($\pm0.06$ Wilson half-width); both headline comparisons clear the noise, and
the fine-tune-vs-full parity is a genuine null rather than an under-powered
one. (ii)~\emph{Single seed}: all runs are seed 0; a second seed of the
low-data pair would strengthen the transfer claim. (iii)~\emph{Full-run
stride}: the full model used T4 stride-2 subsampling, so its $0.23$ is a
conservative lower bound and cross-comparisons against it should note this. (iv)~\emph{Precision}: T4 runs
use fp16, A100 bf16/TF32; results across tiers are statistically equivalent,
not bit-identical. (v)~A log-rotation flaw (since fixed) destroyed the
full run's early-epoch loss lines, so its training curve may begin after epoch
1; the model and all planning results are unaffected.

\section{Addendum: Coupled Local/Global Planning}
\label{sec:localglobal}

The trained DINO-WM checkpoint above is also the \emph{global} model for a
second study: cheap first-order (gradient-based) planning against a trusted,
non-differentiated world model~\cite{fog2026}. Zero-order planners (CEM) work
with any accurate forward model but scale poorly in horizon and action
dimension; first-order planners are sample-efficient per iteration but need a
differentiable model --- and a model small enough to differentiate cheaply is
usually inaccurate enough to be \emph{exploited}, with gradient descent walking
into regions where it is wrong. The coupled local/global recipe decouples the
two roles: the large model produces trajectories and is never differentiated,
while a lightweight \emph{local} surrogate supplies derivatives evaluated
\emph{along the global model's trajectories}. We instantiate it on the
PointMaze DINO-WM model and ask whether first-order refinement through a cheap
surrogate can recover global-CEM-quality plans at a fraction of the
global-model compute. The surrogate is trained in notebook~03 and the planners
evaluated in notebook~08; the full method and protocol are in the companion
file \texttt{reports/local\_global\_report.tex}.

\subsection{Method}
The \textbf{global model} is the frozen DINO-WM latent predictor above, used as
forward simulator and scorer under \texttt{no\_grad}. The \textbf{local model}
is a small residual-MLP surrogate over a frozen orthogonal projection of the
same latents, differentiable in the actions and trained on cached transitions
($20$k steps). Six \textbf{planners} are compared: \texttt{global\_cem}
(zero-order on the global model); \texttt{local\_cem}, \texttt{local\_gd},
\texttt{local\_adam} (zero/first-order on the surrogate alone, isolating model
error from optimizer effects); \texttt{hybrid\_cem\_local\_refine} (global CEM,
then local first-order refinement of the winning sequence); and
\texttt{hybrid\_cem\_local\_refine\_global\_rescore} (the same, but the global
model re-scores the refinement and keeps it only if it improves the global
objective). Three safeguards against surrogate exploitation are built in --- the
frozen orthogonal projection (anti-collapse), an action-smoothness penalty, and
the global re-score gate --- and the local--global \emph{disagreement} (surrogate
prediction vs.\ the global model's executed step) is logged throughout.

\subsection{Protocol and its stated bias}
Planners are scored on offline latent goal reaching over $100$ held-out
episodes: normalized final distance
$\rho=\lVert\zlat_{\mathrm{fin}}-\zlat_{\mathrm{goal}}\rVert^2/
\lVert\zlat_{\mathrm{start}}-\zlat_{\mathrm{goal}}\rVert^2$ and
$\text{success}=\mathbb{1}[\rho<0.5]$. \textbf{Because the global model is both
simulator and judge}, \texttt{global\_cem} is structurally favored on distance
--- it optimizes exactly the quantity it is scored on --- so the
distance/success columns establish only that the other planners are
\emph{competent}, while the efficiency columns (wall time, forward/backward
call counts) and the acceptance rate carry the comparative claim. A
\emph{reference replay} of the dataset's true action sequence through the same
simulator calibrates the global model's own rollout error at
$\rho_{\mathrm{ref}}=0.058$.

\subsection{Surrogate training}
The surrogate reaches a final validation one-step loss of $0.0058$ and a
$K$-step rollout MSE of $0.0029$ at $20{,}000$ steps (per-horizon curve in
Figure~\ref{fig:lg-rollout}).

\begin{figure}[h]
\centering
\drivefigure{pointmaze_local_full_seed0_rollout_error.png}{0.6\linewidth}{pointmaze\_local\_full\_seed0\_rollout\_error.png}
\caption{Surrogate validation rollout error vs.\ horizon (PointMaze, seed~0),
exported by notebook~08 to the Drive \texttt{\detokenize{_summary}} directory.}
\label{fig:lg-rollout}
\end{figure}

\subsection{Planner comparison}
\begin{table}[h]
\centering
\scriptsize
\setlength{\tabcolsep}{4pt}
\begin{tabular}{@{}lcccccccc@{}}
\toprule
Planner & succ. & $\rho$ & wall/ep.\ (s) & glob.\ fwd & loc.\ fwd & bwd & accept & disagr. \\
\midrule
\texttt{global\_cem}                     & 0.97 & 0.138 & 231.6 & 1.50M & --    & --    & --    & --     \\
\texttt{local\_cem}                      & 0.45 & 0.675 & 0.22  & --    & 1.50M & --    & --    & 0.0020 \\
\texttt{local\_gd}                       & 0.13 & 1.306 & 2.60  & --    & 50k   & 50k   & --    & 0.0028 \\
\texttt{local\_adam}                     & 0.41 & 0.699 & 2.67  & --    & 50k   & 50k   & --    & 0.0019 \\
\texttt{hybrid\_cem\_local\_refine}      & 0.54 & 0.632 & 234.1 & 1.50M & 50k   & 50k   & 1.00  & 0.0019 \\
\texttt{hybrid\_\dots\_global\_rescore}  & \textbf{0.97} & \textbf{0.138} & 234.2 & 1.50M & 50k & 50k & 0.028 & 0.0024 \\
\midrule
reference replay                         & --   & 0.058 & --    & --    & --    & --    & --    & --     \\
\bottomrule
\end{tabular}
\caption{Planner comparison on offline latent goal reaching, $100$ held-out
PointMaze episodes (\texttt{pointmaze\_local\_full\_seed0}, surrogate at
$20$k steps). Distances are judged by the global model (see the stated bias),
so the efficiency columns carry the comparative claim. \texttt{succ.}=success
rate; $\rho$=normalized final distance; \texttt{fwd/bwd}=forward/backward call
counts; \texttt{accept}=refinement acceptance rate;
\texttt{disagr.}=mean local--global disagreement.}
\label{tab:lg-planners}
\end{table}

\subsection{Findings}
\begin{enumerate}[leftmargin=1.5em]
\item \textbf{Global CEM tops the distance metric} ($0.97$ success,
$\rho=0.138$), as the judge bias predicts; even so, its $0.138$ is
${\approx}2.4\times$ the reference-replay floor ($0.058$) --- the global model
does not perfectly close its own gap.
\item \textbf{Pure local planners are far cheaper but weaker} on the
global-judged metric: \texttt{local\_cem} $0.45$ at \textbf{$0.22$\,s/episode},
\texttt{local\_adam} $0.41$ at $2.7$\,s, \texttt{local\_gd} $0.13$ --- against
\texttt{global\_cem}'s \textbf{$231$\,s/episode}. The surrogate forward is
${\sim}1000\times$ cheaper per call, but at this capacity it is not accurate
enough to plan to global quality on its own.
\item \textbf{The unguarded hybrid hurts.}
\texttt{hybrid\_cem\_local\_refine} accepts \emph{every} refinement
($\text{accept}{=}1.00$) and falls to $0.54$ ($\rho=0.632$): refining on the
surrogate walks a good global plan into regions where the surrogate is wrong
--- exactly the exploitation failure mode the safeguards target.
\item \textbf{The global re-score gate fixes it.}
\texttt{hybrid\_\dots\_global\_rescore} rejects $97.2\%$ of refinements
($\text{accept}{=}0.028$) and recovers global-CEM quality \emph{exactly}
($0.97$, $\rho=0.138$). This is the headline result: the gate turns a harmful
refinement into a no-op, so the hybrid is never worse than the CEM it began
from.
\item \textbf{Disagreement is small, yet gradients are unreliable.} The
local--global disagreement stays ${\approx}0.002$ --- the surrogate tracks the
global model's one-step behavior --- yet its gradient \emph{directions} improve
the global objective rarely enough that the gate accepts only $2.8\%$ of them.
One-step tracking accuracy and gradient usefulness are not the same thing.
\end{enumerate}

\subsection{Synthetic sanity check (preliminary)}
On the exact synthetic point-mass task (a few hundred surrogate steps; 2--3
episodes, a correctness signal rather than statistics) the expected ordering
holds: all six planners succeed, \texttt{local\_adam} attains the lowest final
distance, the hybrids improve on plain CEM with their refinements accepted, and
\texttt{local\_gd} is weakest. In an adversarial unit test where the surrogate's
action gain is sign-flipped, refinement reliably worsens the global cost and the
re-score gate rejects it, returning the CEM sequence unchanged.

\subsection{What this establishes}
The protocol measures planner behavior \emph{under the global model's belief}.
It shows that first-order refinement through a cheap surrogate can be made
\emph{safe} --- never worse than CEM --- by a global re-score gate, and that the
pure local planners are $100$--$1000\times$ cheaper per episode than the global
baseline. It cannot, by construction, rank planners against reality;
real-environment confirmation in MuJoCo PointMaze is the headline next step
(\S\ref{sec:future}). Results are single-seed (\texttt{seed0}), single-task.

\section{Computational resources}
\label{sec:resources}

A standing constraint of this project is that every experiment fit inside
free-tier Colab limits, and that the experiment \emph{definition} (model size,
data budgets, episode counts) be identical across GPU tiers --- only throughput
knobs (precision, dataloader workers, micro-batch size) differ between a T4, an
L4, and an A100 (\S\ref{sec:implementation}). The figures below therefore
describe a pipeline that runs end-to-end on free hardware, not a cluster job.

\paragraph{Hardware.} Baseline runs use a free-tier NVIDIA T4 (16\,GB, fp16); an
A100-40\,GB or L4 (bf16) is used opportunistically when Colab grants one. Turing
emulates bf16 slowly, so the T4 runs in fp16 while A100/L4 use bf16; results
across tiers are statistically equivalent, not bit-identical.

\paragraph{The enabling optimization.} The stock DINO-WM loop re-encodes
${\sim}1600\times$ the dataset through the frozen DINOv2 encoder every epoch,
projecting $19$--$30$\,h per epoch on a T4 and putting the full study out of
reach on free hardware. Precomputing the frozen patch latents once and training
the predictor on that cache cuts per-step cost ${\approx}21\times$
(Table~\ref{tab:throughput}): a measured $31$\,min/epoch on a T4 at stride~2, or
$3$--$6$\,min on an A100. The cache is ${\sim}30$\,GB for all $2{,}200$ episodes
and lives on session-local disk.

\begin{table}[h]
\centering
\small
\begin{tabular}{@{}lll@{}}
\toprule
Stage & Budget & Wall-clock \\
\midrule
Latent precompute (one-time)         & $2{,}200$ episodes $\to$ DINOv2      & amortized; ${\sim}30$\,GB cache \\
Predictor training, full             & $10$ epochs, $2{,}000{+}200$ traj.  & ${\sim}31$\,min/epoch T4; $3$--$6$\,min A100 \\
Predictor training, low-data ($\times2$) & $20$--$30$ epochs, $300{+}60$ traj. & well under the full run ($15\%$ data) \\
Planning eval (per checkpoint)       & $200$ episodes, CEM $100{\times}10$  & ${\sim}47$\,min (A100, measured) \\
Surrogate training (\S\ref{sec:localglobal}) & $20{,}000$ steps             & ${\sim}1$--$2.5$\,h (T4) \\
Six-planner eval (\S\ref{sec:localglobal})   & $100$ episodes $\times\,6$ planners & ${\sim}20$\,h total (see below) \\
\bottomrule
\end{tabular}
\caption{Wall-clock footprint of both studies on the stated hardware. Training
and planning are \emph{per checkpoint}; three checkpoints (full, low-data
scratch, low-data fine-tune) enter the DINO-WM results.}
\label{tab:resources}
\end{table}

\paragraph{Where the time goes.} The six-planner sweep is dominated by the three
planners that call the global model: each runs ${\approx}232$--$234$\,s per
episode (${\approx}6.5$\,h for $100$ episodes), against ${\approx}0.2$--$2.7$\,s
per episode --- seconds to a few minutes in total --- for the three purely local
planners, a ${\sim}1000\times$ per-call gap. This asymmetry is exactly the
motivation for the coupled local/global design (\S\ref{sec:localglobal}): the
global rollout is the dominant cost, so a cheap differentiable surrogate that
substitutes for some of those rollouts is where any compute saving must come
from.

\paragraph{Resumability.} Because several stages exceed a single free-tier
session --- the six-planner sweep alone is ${\sim}20$\,h --- every long-running
stage is checkpointed and resumable: training writes per-epoch checkpoints,
planning resumes per episode (the latest \texttt{final\_eval} is authoritative),
and the local/global evaluation runs under a $300$-minute per-session wall cap
with cross-session resume. No stage needs to finish in one sitting, which is
what makes a ${\sim}20$\,h evaluation tractable on free, time-limited GPUs.

\section{Future work}
\label{sec:future}

(1)~Real-environment confirmation of the local/global planner ranking in MuJoCo
PointMaze (the offline metric of \S\ref{sec:localglobal} judges distances with
the global model itself); (2)~seed replication of the low-data pair and the
local/global run; (3)~an A100
stride-1 full-data training to replace the stride-2 lower bound;
(4)~PushT, blocked upstream on a video-backed latent cache (the PointMaze
latent path does not transfer directly); and (5)~for the local/global track,
FOG-style online surrogate fine-tuning on global-model rollouts to keep local
gradients on-distribution, plus the built-in projection/context/optimizer
ablations --- both motivated by the low refinement-acceptance rate of
\S\ref{sec:localglobal}.

\section*{Reproducibility}

Code lives in the \texttt{wm-prediction} repository
(\texttt{src/wm\_poc/dino\_wm}, \texttt{scripts/dino\_wm},
\texttt{configs/dino\_wm}); notebook 02 runs the full pipeline on Colab
(T4 or A100) with self-gating, resumable stages, and notebook 07 renders the
results. The latent-cache design is documented in
\texttt{docs/dino\_wm\_latent\_cache\_training.md}; the figures in this report
are generated from the recorded planner telemetry
(\texttt{reports/pointmaze\_planning\_logs\_50evals.csv}) by
\texttt{reports/make\_planning\_figures.py}. The two training-curve figures are
produced on Drive by notebook 07 (\texttt{\_summary/}) and copied into
\texttt{reports/figures/}.

\begin{thebibliography}{9}
\setlength{\itemsep}{2pt}

\bibitem{dinowm2025}
G.~Zhou, H.~Pan, Y.~LeCun, L.~Pinto.
\emph{DINO-WM: World Models on Pre-trained Visual Features Enable Zero-shot
Planning}. arXiv:2411.04983, 2024.

\bibitem{dinov2}
M.~Oquab et al.
\emph{DINOv2: Learning Robust Visual Features without Supervision}.
arXiv:2304.07193, 2023.

\bibitem{planet}
D.~Hafner et al.
\emph{Learning Latent Dynamics for Planning from Pixels}. ICML, 2019.

\bibitem{dreamer}
D.~Hafner et al.
\emph{Dream to Control: Learning Behaviors by Latent Imagination}. ICLR, 2020.

\bibitem{tdmpc2}
N.~Hansen, H.~Su, X.~Wang.
\emph{TD-MPC2: Scalable, Robust World Models for Continuous Control}.
ICLR, 2024.

\bibitem{vjepa2}
M.~Assran et al.
\emph{V-JEPA 2: Self-Supervised Video Models Enable Understanding, Prediction
and Planning}. 2025.

\bibitem{fog2026}
P.~Amigo, R.~Khorrambakht, N.~Mansard, L.~Righetti.
\emph{Coupled Local and Global World Models for Efficient First-Order
Model-Based Reinforcement Learning}. arXiv:2602.06219, 2026.

\bibitem{cem}
P.-T.~de Boer, D.~P.~Kroese, S.~Mannor, R.~Y.~Rubinstein.
\emph{A Tutorial on the Cross-Entropy Method}. Annals of Operations Research,
134(1):19--67, 2005.

\end{thebibliography}

\end{document}
